Some implications and consequences of the expansion of the universe are examined. 
In Chapter 1 it is shown that this expansion creates grave difficulties for the Hoyle–Narlikar theory of gravitation. 

Chapter 2 deals with perturbations of an expanding homogeneous and isotropic universe. 
The conclusion is reached that galaxies cannot be formed as a result of the growth of perturbations that were initially small. 
The propagation and absorption of gravitational radiation is also investigated in this approximation. 

In Chapter 3 gravitational radiation in an expanding universe is examined by a method of asymptotic expansions. 
The ‘peeling–off’ behaviour and the asymptotic group are derived. 

Chapter 4 deals with the occurrence of singularities in cosmological models. 
It is shown that a singularity is inevitable provided that certain very general conditions are satisfied.
